\documentclass[11pt,letterpaper]{article}
\usepackage[margin=1in]{geometry}
\usepackage{amsmath,amssymb}
\usepackage{booktabs,array}
\usepackage{enumitem}
\usepackage{xcolor}
\usepackage{microtype}
\usepackage{tcolorbox}
\tcbuselibrary{breakable,skins}

\newtcolorbox{auditpass}{breakable,enhanced,
  colback=green!3,colframe=green!55!black,
  fonttitle=\bfseries\small,title=Audited: Correct and New,
  boxrule=0.8pt,left=4pt,right=4pt,top=3pt,bottom=3pt}
\newtcolorbox{auditfix}{breakable,enhanced,
  colback=orange!3,colframe=orange!65!black,
  fonttitle=\bfseries\small,title=Audited: Contains Error — Corrected Below,
  boxrule=0.8pt,left=4pt,right=4pt,top=3pt,bottom=3pt}
\newtcolorbox{keynum}{breakable,enhanced,
  colback=blue!3,colframe=blue!55!black,
  fonttitle=\bfseries\small,title=Key Number,
  boxrule=0.8pt,left=4pt,right=4pt,top=3pt,bottom=3pt}

\newcommand{\kB}{k_{\mathrm{B}}}
\newcommand{\asc}{\alpha_{\mathrm{screen}}}
\newcommand{\Scert}{S_{\mathrm{cert}}}
\newcommand{\Leff}{\Lambda_{\mathrm{eff}}}

\title{\textbf{Phase III--IV Integration: Control Architecture,\\
CP Shielding, and Deployment Roadmap}\\[4pt]
\large Audited Engineering Modules for Project SNR--Quantum Gravity\\[3pt]
\normalsize Supplement to: \textit{Entanglement-Entropy Sourcing of Gravity}, v4}
\author{Kevin Monette}
\date{March 2026}

\begin{document}
\maketitle

This document integrates the genuinely new engineering contributions from the
``ship-ready blueprint'' into the project's canonical framework, with each claim
independently audited.  We incorporate what is correct, correct what has an error,
and discard what duplicates material already in the project.

% ─────────────────────────────────────────────────────────────────────────────
\section{Revised Integration Time: 34 Seconds, Not 17}

\begin{auditpass}
The revised integration time of 34\,s is correct and represents a genuine
upgrade over the 17\,s figure in the previous addendum.
\end{auditpass}

The 17\,s calculation in the Strategic Addendum used only the statistical noise floor
$\sigma_{\mathrm{noise}} = 3\times10^{-9}$\,m/s$^2$/\!\sqrt{Hz}.  When the residual
Casimir-Polder systematic $\sigma_{\mathrm{sys}} \approx 0.5\times10^{-9}$\,m/s$^2$
(from 0.05\% shield non-uniformity, treated after calibration) is included:
\begin{equation}
  \sigma_{\mathrm{tot}}^2 = \frac{\sigma_{\mathrm{noise}}^2}{T_{\mathrm{int}}}
  + \sigma_{\mathrm{sys}}^2.
\end{equation}
Setting Signal\,$/\sigma_{\mathrm{tot}} = 5$ with
Signal $= 3.6\times10^{-9}$\,m/s$^2$:
\begin{align}
  \left(\frac{3.6\times10^{-9}}{5}\right)^2
  &= \frac{(3\times10^{-9})^2}{T_{\mathrm{int}}} + (0.5\times10^{-9})^2, \\
  T_{\mathrm{int}}
  &= \frac{(3\times10^{-9})^2}{(0.72\times10^{-9})^2 - (0.5\times10^{-9})^2}
  = \frac{9\times10^{-18}}{(0.52-0.25)\times10^{-18}}
  \approx 33.5\,\text{s}.
\end{align}

\begin{keynum}
Revised 5$\sigma$ integration time: $T_{\mathrm{int}} \approx \mathbf{34}$\,\textbf{s},
accounting for Casimir-Polder systematic after calibration subtraction.
Per-run data requirement: 100 valid runs $\times$ 34\,s $\approx$ 57\,minutes total.
\end{keynum}

The 34\,s figure replaces the 17\,s figure in all subsequent project documents.

% ─────────────────────────────────────────────────────────────────────────────
\section{Gold-on-Silicon Faraday Shield for CP Mitigation}

\begin{auditpass}
The 200\,nm Gold-on-Silicon Faraday shield is a genuine new component, correctly
designed to suppress Casimir-Polder forces while being gravitationally transparent.
\end{auditpass}

\subsection{Design specification}

\begin{center}\small
\begin{tabular}{@{}llll@{}}
\toprule
Parameter & Specification & Basis & Status \\
\midrule
Material & 200\,nm Au on Si substrate & High conductivity; mechanical stability & Manufacturable \\
Geometry & Grounded planar shield & Screens static $E$-field; suppresses CP & Standard \\
Gravitational transparency & $<10^{-12}$ attenuation & Gravity couples to $T_{\mu\nu}$; 200\,nm Au is negligible & Trivially satisfied \\
Thickness control & $\pm100$\,nm (0.05\%) & ALD deposition; AFM verification & Achievable \\
\bottomrule
\end{tabular}
\end{center}

\subsection{Residual systematic from non-uniformity}

A 100\,nm peak-to-peak thickness variation at $R = 100\,\mu$m produces a Casimir
pressure gradient calculable via the Lifshitz formula.  The estimated systematic
acceleration on the SQUID lever (1\,$\mu$g, 1\,$\mu$m$^2$ tip) after topographic
mapping and calibration subtraction is $\sigma_{\mathrm{sys}} \approx 0.5\times10^{-9}$\,m/s$^2$.
This number enters the 34\,s integration time above.

\textbf{Mitigation path:}  AFM mapping of the fabricated shield provides a topographic
model $\delta h(x,y)$.  The corresponding Casimir gradient is computed from the Lifshitz
formula and subtracted as a prior in the Bayesian analysis (Section~\ref{sec:bayes}).

% ─────────────────────────────────────────────────────────────────────────────
\section{MPC Thermal-Nulling Control Loop}

\begin{auditpass}
The Model Predictive Control (MPC) architecture with look-ahead from the gate schedule
is a genuine conceptual improvement over a standard PID controller.  The key idea ---
using the scheduled gate sequence to anticipate heat deposition before it occurs ---
is not in any prior project document.
\end{auditpass}

\begin{auditfix}
The blueprint states the thermal diffusion time across a 500\,$\mu$m silicon substrate
is $\tau_{\mathrm{th}} \approx 25$\,ns, using $D = 10$\,m$^2$/s.  This is the room-temperature
diffusivity.  At $T = 15$\,mK, silicon has dramatically reduced thermal conductivity
and heat capacity.  The correct values are $\kappa_{\mathrm{Si}} \approx 0.01$\,W/(m\,K)
and $c_p \approx 10^{-10}$\,J/(kg\,K) (Debye $T^3$ suppression), giving
$D_{\mathrm{mK}} \approx 50{,}000$\,m$^2$/s and $\tau_{\mathrm{th}} \approx 5$\,ps.
Heat in fact spreads \emph{faster} at 15\,mK through ballistic phonons; the thermal
equilibration of the chip is near-instantaneous on any timescale relevant to the experiment.
The MPC latency requirement of $\tau_{\mathrm{loop}} \leq 5\,\mu$s is driven by the
\emph{electrical} control system, not thermal diffusion.  The thermal-diffusion argument
in the blueprint is wrong, but the $5\,\mu$s latency requirement is correct on independent
grounds.
\end{auditfix}

\subsection{Corrected MPC architecture}

The MPC controller operates as follows:
\begin{enumerate}[noitemsep]
  \item The quantum pulse generator provides a look-ahead vector of expected energy
    deposition $\vec{E}_{\mathrm{gate}}(t)$ for the next $\Delta t = 5\,\mu$s window.
  \item TES bolometers (NEP $\approx 10^{-20}$\,W/$\sqrt{\mathrm{Hz}}$) at the thermal
    anchors of both branches read $P_A(t)$ and $P_B(t)$ with power resolution
    $\approx 4.5\times10^{-18}$\,W in a 5\,$\mu$s window.
  \item The MPC optimizer minimizes:
    \begin{equation}
      J = \sum_{k=0}^{N_p}\!\left(|\Delta P_{\mathrm{diss}}(t+k)|^2
          + \lambda|\vec{E}_{\mathrm{comp}}(t+k)|^2\right)
    \end{equation}
    and emits compensation tones at frequency-detuned frequencies to dummy resistive
    elements, matching the heat deposited in branch A to branch B.
  \item The FPGA implementation achieves loop latency $\tau_{\mathrm{loop}} \leq 5\,\mu$s
    using current commercial hardware (e.g., Keysight M3300A AWG, Xilinx Ultrascale+).
\end{enumerate}

\textbf{Why this is better than PID:}  A standard PID controller reacts to errors
\emph{after} they occur.  Heat dissipated by a CZ gate cannot be ``un-deposited.''
The MPC look-ahead computes the compensation pulse \emph{before} the gate fires,
preventing the imbalance rather than correcting it.  This is the correct architecture
for matching dissipation to the required $10^{-16}$\,W precision.

% ─────────────────────────────────────────────────────────────────────────────
\section{Bayesian Signal Extraction and Blind Injection}
\label{sec:bayes}

\begin{auditpass}
The Bayesian hypothesis-testing framework and blind injection protocol are both
genuine additions.  The Bayes factor formulation correctly handles the systematic floor.
The blind injection requirement (identify $>90\%$ of injections, $<5\%$ false positives)
is a standard good-practice requirement not previously specified.
\end{auditpass}

\subsection{Statistical framework}

The measured lever displacement stream $d(t)$ is analyzed under two hypotheses:
\begin{align}
  H_0:&\quad d \sim \mathcal{N}(0,\, \sigma_{\mathrm{noise}}^2 + \sigma_{\mathrm{CP}}^2), \\
  H_1:&\quad d \sim \mathcal{N}(\Delta a,\, \sigma_{\mathrm{noise}}^2 + \sigma_{\mathrm{CP}}^2).
\end{align}
The Bayes factor is $K = P(d|H_1)/P(d|H_0)$.  A detection is claimed at
$\ln K > 5$ (approximately $5\sigma$ in the Gaussian limit).

With $\Delta a = 3.6\times10^{-9}$\,m/s$^2$, $\sigma_{\mathrm{noise}}
= 3\times10^{-9}$\,m/s$^2/\sqrt{\mathrm{Hz}}$, and
$\sigma_{\mathrm{CP}} = 0.5\times10^{-9}$\,m/s$^2$ (Lifshitz-model systematic):
\begin{equation}
  T_{\mathrm{int}}(\text{Bayesian}) = 33.5\,\text{s},
\end{equation}
consistent with the direct calculation above.

\subsection{Blind injection protocol}

\begin{itemize}[noitemsep]
  \item 10\% of data runs receive a synthetic signal injected digitally into the SQUID
    readout by a trusted third-party algorithm, without disclosure to the analysis team.
  \item Success criterion: $>90\%$ of injections identified, $<5\%$ false positives.
  \item Pipeline validated against blind injections before unblinding actual data.
\end{itemize}

This protocol prevents analyst bias and validates the entire signal extraction pipeline.

% ─────────────────────────────────────────────────────────────────────────────
\section{18-Month Deployment Roadmap}

\begin{center}\small
\begin{tabular}{@{}p{2cm}p{4cm}p{5cm}p{2.5cm}@{}}
\toprule
Phase & Duration & Deliverable & Milestone \\
\midrule
\textbf{Theory} & Now & Submit Paper I (v4); OP3 Stage 1 code & arXiv submission \\
\addlinespace
\textbf{Fabrication} & Months 1--6 & ALD Au/Si shields; 50-qubit transmon array with integrated TES heaters & AFM: shield uniformity $<100$\,nm ripple \\
\addlinespace
\textbf{Calibration} & Months 7--9 & Cool to 15\,mK; characterize $T_2$; tune MPC loop & $|\Delta P_{\mathrm{diss}}| \leq 10^{-16}$\,W stable for 1\,h \\
\addlinespace
\textbf{Validation} & Months 10--12 & Blind injection protocol; baseline noise ($H_0$); OP3 Stage 3 complete & Injections identified; $\ln K < 5$ for baseline \\
\addlinespace
\textbf{Execution} & Months 13--15 & Data collection, squeezed cluster state, 100 runs $\times$ 34\,s & Accumulated $3{,}400$\,s valid data \\
\addlinespace
\textbf{Publication} & Months 16--18 & Unblind; compute Bayes factor; draft paper & Publish (positive or null) \\
\bottomrule
\end{tabular}
\end{center}

% ─────────────────────────────────────────────────────────────────────────────
\section{Revised System Integration Matrix}

\begin{center}\small
\begin{tabular}{@{}p{2.5cm}p{2.5cm}p{2.5cm}p{3.5cm}p{2cm}@{}}
\toprule
Subsystem & Module & Interface & Critical constraint & Status \\
\midrule
Quantum source & MPO-TDVP (OP3) & Pulse generator & $\chi \leq 256$; $\asc \geq 10^{-2}$ & Pending Stage 3 \\
CP shield & Faraday (Au/Si) & Mechanical mount & Non-uniformity $< 0.05\%$ & Design complete \\
Thermal control & MPC with TES & FPGA + bolometers & Latency $\leq 5\,\mu$s; $\Delta P \leq 10^{-16}$\,W & Architecture set \\
Data analysis & Bayesian ($H_0/H_1$) & SQUID readout & $T_{\mathrm{int}} \geq 34$\,s; blind injection pass & Framework complete \\
\bottomrule
\end{tabular}
\end{center}

% ─────────────────────────────────────────────────────────────────────────────
\section{Updated Risk Register}

\begin{center}\small
\begin{tabular}{@{}p{3cm}p{1.5cm}p{1.5cm}p{4.5cm}p{2.5cm}@{}}
\toprule
Risk & Prob. & Impact & Mitigation & Gates \\
\midrule
Bond dimension $\chi > 256$ & Medium & High & Reduce to $N=30$; longer integration & Blocks Paper III \\
Shield non-uniformity $>0.05\%$ & High & Medium & AFM map; Lifshitz prior in Bayes & Limits SNR \\
MPC latency $> 5\,\mu$s & Low & High & FPGA pre-computation; look-ahead & Required \\
$\asc(50,T_\phi) < 10^{-4}$ & Unknown & Critical & Revert to $N=30$ or revise $E_*$ hypothesis & Pivotal \\
CP calibration fails & Low & Medium & Increase shield uniformity spec & Increases $T_{\mathrm{int}}$ \\
\bottomrule
\end{tabular}
\end{center}

% ─────────────────────────────────────────────────────────────────────────────
\section{Feasibility Assessment}

The blueprint's overall GREEN/``ship-ready'' conclusion is premature.  The correct
current status is:

\begin{center}\small
\begin{tabular}{@{}p{4cm}p{3cm}p{5.5cm}@{}}
\toprule
Module & Status & Blocking issue \\
\midrule
Theory (v4 paper) & \textbf{Ready to submit} & None \\
MPO-TDVP (OP3 Stage 1) & Not started & Requires 2--4 weeks QuTiP code \\
MPO-TDVP (OP3 Stage 3) & Not started & Requires 4--6 months tensor network work \\
Faraday shield design & Complete & Fabrication requires facility \\
MPC architecture & Specified & FPGA implementation: 2--3 months \\
Bayesian pipeline & Framework set & Requires OP3 output for priors \\
Full experiment & \textbf{Not yet authorized} & Requires OP3 Stage 3 + $\Delta P$ calorimetry \\
\bottomrule
\end{tabular}
\end{center}

The experiment is not ``ready to fabricate'' in the sense of authorizing all phases.
It is ready to:
(1)~submit Paper I immediately,
(2)~begin OP3 Stage 1 numerics in parallel,
(3)~engage a fabrication facility for shield samples and test bolometers.

The full hardware integration should be authorized only after the OP3 go/no-go
criterion ($\asc(50,T_\phi) \geq 10^{-2}$) is met.

\end{document}
